\documentclass[uplatex, dvipdfmx]{jsarticle}

\include{begin}

\section{ラジアン}

ラジアンとは、半径と同じ長さの弧を切り取る2本の半径が成す角の値。

したがって、1ラジアンの角度は、以下の計算で求められる。

\[
 1radian
 = 360 \times \frac{半径}{円周}
 = 360 \times \frac{r}{2\pi r}
 = \frac{180}{\pi}
\]

\section{32度は何ラジアンか？}

20m離れたところにある木の高さは、

角度32度のラジアン値を x とすると、

距離20m $\times$ tan(x) によって求めることができる。

\[
 xラジアン : 1ラジアン = 32度 : \frac{180}{\pi}
\]
\[
 x = 1 \times 32 \div \frac{180}{\pi}
\]
\[
 x = 32 \times \frac{\pi}{180}
\]

\include{end}

%% 修正時刻： Mon 2024/10/07 20:19:000
